\newpage
\section{Installation}\label{installation} 

The current version of the extension can be downloaded from
\url{https://bitbucket.org/Klinkenstecker/csurf\_klinke\_mcu/downloads}.
To install the extension:
\begin{enumerate}
\item open the \reaper resources path, e.g. via
\reaper by selecting the menu-entry: {\tt Options/ Show \reaper
  resource path in explorer/finder...}.
\item In this path, there is a
  folder with the name Locate the {\tt UserPlugin} folder within the \reaper resource path.
\item Unpack the downloaded
  zip archiv into the {\tt UserPlugin} directory.
\end{enumerate}

  Afterwards the {\tt UserPlugin} folder
should contains four new {\tt reaper\_csurf\_mcu\_klinke} dll-files,
this manual and a new folder named {\tt MCU}. This {\tt MCU} folder
contains default configuration files and FX maps. Do not modify
these files. Your own configuration files and FX maps will be stored in the
folder {\tt My Documents\textbackslash Reaper\textbackslash MCU}. This folder
is created when you start the \mcu surface for the first time.\\

\noindent
To activate the extension you need to follow these steps:
\begin{compactitem}
\item open the \reaper preferences
\item scroll down and select ``Control Surfaces'' 
\item if the original ``Mackie Control Universal'' surface was added, remove it
\item add the ``Mackie Control Protocol (Klinke)'' surface
\item in the case that the controller is connected via USB: increase
  the ``Control surface display update frequency'' in the Control
  surfaces preferences of \reaper to 30 Hz.  
\end{compactitem}

%If you have extenders:
%\begin{itemize}
%\item set the Bank Size parameter in the "Mackie Control Universal (Klinke)"
%preferences dialog to the total number of channels (or maybe to 8 less, in the
%case that you want to use the main unit in non-mixing modes the most time)
%\item remove original "Extender" surfaces and add "Extentder (Klinke)" surfaces
%for each of your extenders
%\item set the Surface Offset parameter in the Surface prefercences dialogs
%corresponding to the order of your setup (the leftmost unit should have an
%offset of 0, the unit on his right side an offset of 8 ...)
%\end{itemize}

% \noindent
% For the \actionmode it is useful to import Action
% shortcuts:\footnote{If you want to assign all actions by yourself
%   there is need to import anything, but  the predefined actions can be
%   useful to understand the features described in the
%   \hyperref[actionmode]{Action-Mode} section.}
 
% \begin{compactitem}
% \item open the \hyperref[F:actiondialog]{Actions dialog} in
%   \reaper(e.g. by pressing the '?' key)
% \item press the Import/Export button on the lower left corner and select
%   "Import..."
% \item in the appearing file-dialog, select the file {\tt \ldots\tbs
%     Reaper\tbs Plugins\tbs MCU}\newline{\tt \tbs
%     DefaultActions.reaperKeyMap}.
% \end{compactitem}

\subsection{Control Surface Settings}\label{settings}

The Control Surface Settings dialog configures the hardware units of
the surface and a few global options.

\subsubsection{Units}\label{units}
The extension combines up to eight MCU-protocol hardware units into
one logical surface: one main unit plus up to seven extenders, i.e.\
up to 64 channels controlled at the same time. Select the unit you
want to edit (Unit~1\,--\,8) and configure it with:
\begin{itemize}
\item \textbf{Type}: the hardware model and role of the unit ---
  \textsc{Mackie Main}, \textsc{Mackie Extender}, \textsc{QCon ProX},
  \textsc{QCon ProX Extender} or \textsc{Disabled}. The two
  \textsc{QCon ProX} types activate the additional hardware elements
  of the Icon QCon Pro~X (second display, master meterbridge) for that
  unit; see section~\ref{prox}.
\item \textbf{MIDI input} / \textbf{MIDI output}: the MIDI ports the
  unit is connected to.
\end{itemize}
The unit position (Unit~1\,--\,8) determines the channel-strip order
from left to right: Unit~1 controls channels~1\,--\,8, Unit~2
channels~9\,--\,16, and so on. Unit~1 can not be disabled, and the
active units must be contiguous from Unit~1 (no gaps in between). A
``main'' unit --- one with transport controls and the
assignment/timecode displays --- may sit at any position; the
transport LEDs and the timecode/assignment displays are sent to every
main unit. The device type fully determines a unit's protocol
behaviour, so there is no separate device-id or offset to set.

\subsubsection{Options}
The following global options can be activated:
\begin{itemize}
\item Use keyboard modifier: The Behringer BCF does not have
  \option/\control and \alt modifier buttons. When this option is
  selected, the keyboard modifier can be used instead. The left
  \alt key substitutes the \alt button, the right \alt key the
  \option button.
\item Fake fader touch: The BCF2000 does not have fader touch sensitive
  faders, so after a fader movement the extension simulates that the
  fader being touched for an additional 2 seconds.
\item Swap left/right arrow on zoom: By default the left arrow key
  zooms out of the arrangement and the right arrow button zooms in (assuming
  zooming is activated via the \zoom button). This can be swapped.
\end{itemize}

\noindent
The former ``2nd controller (Mackie Control B)'' workaround is no
longer needed: a second and further hardware units are now added as
additional units of the same surface, as described above.
%%% Local Variables: 
%%% mode: latex
%%% TeX-master: "../mcu_klinke_manual"
%%% End: 
